\chapter{综合练习}
\section{综合例题}
\noindent\textbf{要点与方法.}\enspace
综合题的计算对象可以重新选择。随机快速排序的比较次数可拆成各对键是否比较的指示量之和，求期望时不要求这些指示量独立；非对称损失的最优点可由凸性和左右导数确定，分布有原子时须保留两侧条件。含无限求和的优化题则应先在有限截断上建立不等式，再说明极限为何成立。

求出候选答案后，还要分别证明它可行、达到界限以及何时取等。平均性能不能保证最坏情形表现，最优解也未必唯一。
\par\smallskip


\problem{CX01}{随机快速排序的比较次数}{}
对 $n$ 个互异键使用随机快速排序：每个非空子问题独立均匀选择一个主元，与该子问题的其余每个键比较一次，再递归处理两侧。设总比较次数为 $C_n$。求 $\E C_n$ 的精确公式，并解释为什么平均为 $O(n\log n)$ 不代表每次运行都是这个量级。
\solution 按键值大小将键编号为 $1,\ldots,n$。令 $I_{ij}$ 表示键 $i<j$ 在算法中是否比较过。每次比较包含一个主元，该主元随后不再进入子问题，所以一对键至多比较一次，并且
\[
C_n=\sum_{1\le i<j\le n}I_{ij}.
\]
只观察秩区间 $\{i,i+1,\ldots,j\}$ 中最先被选作主元的键。区间外的主元不会将 $i,j$ 分到两侧；若第一个区间内主元是内部键，它立即分离 $i,j$，之后二者不再比较；若是端点，则两者当时比较。因此
\[
\Prob(I_{ij}=1)=\frac2{j-i+1}.
\]
“第一个均匀”可由等价实现解释：先给所有键赋独立连续随机优先级，每个子问题选择优先级最低者，所得主元选择规律相同；区间内最低优先级在各键间等可能。

线性性不需要各 $I_{ij}$ 独立。按间隔 $d=j-i$ 求和，令 $H_n=\sum_{k=1}^n1/k$，得到
\[
\begin{aligned}
\E C_n
&=2\sum_{d=1}^{n-1}\frac{n-d}{d+1}\\
&=2(n+1)(H_n-1)-2(n-1)\\
&=\boxed{2(n+1)H_n-4n}.
\end{aligned}
\]
公式对 $n=1$ 给出零；$n=0$ 单独定义 $C_0=0$。例如 $n=4$ 时均值为 $29/6$。

若每次主元恰为当前最小值，就会发生 $(n-1)+\cdots+1=n(n-1)/2$ 次比较；该事件具有正概率 $1/n!$。所以期望阶是对随机运行的平均陈述，不是逐次运行的最坏情形保证。
\insight 分析某一对键何时会相遇，比追踪整棵随机递归树更直接。两两指示量可以相关，期望仍可逐项相加。

\problem{CX02}{缺货与积压的最优平衡}{}
非负需求 $D$ 满足 $\E D<\infty$，其分布函数为 $F$。订购量 $q\ge0$ 可为实数，每单位缺货成本为 $c_u>0$，每单位积压成本为 $c_o>0$。最小化
\[
L(q)=c_u\E(D-q)_++c_o\E(q-D)_+.
\]
不假设需求分布连续，给出全部最优订购量。对于 $\Prob(D=0,1,2)=(1/4,1/2,1/4)$，分别求 $(c_u,c_o)=(3,1)$ 与 $(2,1)$ 时的解集。
\solution 每个固定需求下的成本都是 $q$ 的凸分段线性函数，因此 $L$ 也凸。左右差商均被常数 $\max(c_u,c_o)$ 控制，可以在期望内取单侧极限，得到
\[
\begin{aligned}
L'_-(q)&=(c_u+c_o)F(q-)-c_u,\\
L'_+(q)&=(c_u+c_o)F(q)-c_u.
\end{aligned}
\]
对 $q>0$，凸函数最优性的充要条件为左导数不正、右导数不负。因此令 $\tau=c_u/(c_u+c_o)$，全部最优量为
\[
\boxed{\{q\ge0:F(q-)\le\tau\le F(q)\}}.
\]
在边界 $q=0$，只需右导数不负；由于非负需求使 $F(0-)=0\le\tau$，同一公式依然适用。若 $q<0$，形式上延拓的成本随 $q$ 增加严格下降，所以不会遗漏受约束的最优点。

最优量存在：$\tau\in(0,1)$，广义分位数 $\inf\{x:F(x)\ge\tau\}$ 有限并满足上述条件。若分布函数恰在某段保持水平值 $\tau$，整个区间都可以最优；不能未经检查就宣布分位数唯一。

离散例中，$(c_u,c_o)=(3,1)$ 给出 $\tau=3/4$。在 $q=1$ 处 $F$ 跳到 $3/4$，到 $2$ 前维持此值，故全部最优量为 $[1,2]$；两端及其间的期望成本均为 $1$。当 $(c_u,c_o)=(2,1)$ 时 $\tau=2/3$ 严格落在 $q=1$ 的跳跃内，故唯一解为 $q=1$，最小成本为 $3/4$。
\insight 非对称绝对损失选出成本加权分位数。需求均值一般不是这个决策问题的最优答案，原子和分布空档还会影响唯一性。

\problem{CX03}{给定均值的最大熵分布}{}
随机变量 $N$ 取非负整数，且 $\E N=\mu>0$。用自然对数定义熵
\[
H(N)=-\sum_{k=0}^\infty p_k\log p_k,\qquad 0\log0=0.
\]
证明熵有有限上界，求达到最大熵的唯一分布及最大值。允许使用有限分布的对数和不等式，但须处理无限支撑，不得预先假设熵有限。
\solution 令 $r=\mu/(1+\mu)\in(0,1)$，考虑几何概率
\[
q_k=(1-r)r^k,\qquad k\ge0.
\]
它的均值为 $r/(1-r)=\mu$。若 $N$ 服从该分布，则
\[
H(q)=-\log(1-r)-\mu\log r
=(1+\mu)\log(1+\mu)-\mu\log\mu.
\]

为在一般分布下严格给出上界，将状态 $0,\ldots,M-1$ 分开，把 $N\ge M$ 合并成一个尾部状态，记尾概率为 $t_M$。参考分布 $q$ 对该尾部赋概率 $r^M$。有限分布的对数和不等式给出分组熵
\[
H_M\le-\sum_{k<M}p_k\log q_k-t_M\log(r^M).
\]
右端等于
\[
-(1-t_M)\log(1-r)
-\log r\left(\sum_{k<M}kp_k+Mt_M\right).
\]
当 $M\to\infty$，有 $t_M\to0$，而括号内是 $\E\min(N,M)\uparrow\mu$。同时
$H_M=-\sum_{k<M}p_k\log p_k-t_M\log t_M$ 趋于 $H(N)$，即使起初允许后者为无穷也成立，因为 $-t\log t\to0$。取极限得
\[
H(N)\le(1+\mu)\log(1+\mu)-\mu\log\mu<\infty.
\]

还要证明唯一性。以上右端与 $H_M$ 的差是分组相对熵 $D_M\ge0$，其极限等于 $H(q)-H(N)$。若 $p\ne q$，选一个 $j$ 使 $p_j\ne q_j$。对每个 $M>j$，把除 $j$ 外的所有分组再合并，对数和不等式给出
\[
D_M\ge
p_j\log\frac{p_j}{q_j}
+(1-p_j)\log\frac{1-p_j}{1-q_j}>0.
\]
该严格正下界不依赖 $M$，故不可能取等。于是几何分布是唯一最大熵分布。
\insight 在无限状态上，先将尾部合并成一个状态，再取极限，可避免把待证的有限熵假设偷偷用于相对熵运算。


\section{练习题}
\noindent\textbf{要点与方法.}\enspace
综合练习先把条件转成可计算的结构。模约束可用剩余类或单位根筛选；旋转对称性可把条件矩化为球面上的平均，但须注意取样所用的测度；连续试验的模式则可用有限状态记录仍会影响未来的后缀。计算方差时，不能省略共享变量或时间相关产生的协方差项。

整数题常需分离素因子的指数与剩余部分，模运算中的除法只有在分母可逆时才合法。渐近计数应说明边界与重复元素造成的误差。有限枚举可以检查小规模结果，一般结论仍需递推、恒等式或不变量的证明。
\par\smallskip



\exercise{AA01}{局部独立与整体奇偶约束}
独立生成 $n-1$ 个公平比特 $B_1,\ldots,B_{n-1}$，其中 $n\ge3$，再令 $B_n$ 为前面比特之和模2。证明任意 $n-1$ 个比特仍相互独立且公平，但全部 $n$ 个比特不独立。求所有比特均为0的概率。

\exercise{AA02}{成功次数模三的分布}
独立抛掷 $n\ge0$ 枚公平硬币，正面总数为 $S$。求 $\Prob(S\equiv r\pmod3)$，其中 $r=0,1,2$。可使用 $\omega=e^{2\pi i/3}$，但须说明单位根筛选为何成立；最后计算 $n=8,r=0$。

\exercise{AA19}{相邻交换后的复原概率}
三张不同的牌初始排列为123。每一步独立地以相同概率交换位置1与2，或交换位置2与3。求经过 $t\ge0$ 步后恢复123的概率，并解释该概率为何没有单一的长期极限。

\exercise{AA22}{环形排列中的孤立标记}
把 $m$ 个 A 与 $n$ 个 B 随机排列在 $N=m+n\ge3$ 个编号圆周位置，所有类型排列等可能。一个 B 的左右邻居都是 A 时称为孤立 B。求孤立 B 数的期望；再求线性排列中只计非端点位置时的期望。代入 $m=4,n=3$。

\exercise{AA03}{圆盘投影的分布与绝对矩}
在半径 $R>0$ 的圆盘内按面积均匀取点 $(X,Y)$。求横坐标 $X$ 的密度、$\E|X|$ 和 $\Var X$，并解释为什么 $X$ 不均匀分布于 $[-R,R]$。

\exercise{AA07}{圆周两类样本的线性可分性}
在单位圆周上独立均匀选取 $m+n$ 个点，其中前 $m\ge2$ 个标为 A，其余 $n\ge2$ 个标为 B。求存在一条直线把所有 A 与所有 B 严格分在两侧的概率。给出 $m=n=3$ 的结果。

\exercise{AA09}{正象限中的高斯比值}
$X,Y$ 独立标准正态。在条件 $X>0,Y>0$ 下令 $Z=X/Y$。求 $Z$ 的分布函数、密度、中位数，并判断 $\E[Z\mid X>0,Y>0]$ 是否有限。

\exercise{AA10}{给定总能量后的平方相关}
$X,Y,Z$ 独立标准正态，$S=X^2+Y^2+Z^2$。求 $\E[X^2Y^2\mid S]$、$\Var(X^2\mid S)$ 与 $\Cov(X^2,Y^2\mid S)$。请从旋转对称性推导所需四阶矩，不直接引用球面矩公式。

\exercise{AA13}{三个到达时刻的最近间隔}
三个到达时刻独立均匀分布于 $[0,L]$，其中 $L>0$。记任意两次到达之间的最短时间差为 $D$。求 $D$ 的完整分布函数和期望，不把到区间端点的距离计作间隔。

\exercise{AA16}{指数样本的最小值与极差}
$X_1,\ldots,X_n$ 独立服从参数 $\lambda>0$ 的指数分布，$n\ge2$。令 $U=\min_iX_i$、$V=\max_iX_i$、$D=V-U$。证明 $U,D$ 独立，求 $D$ 的分布函数及 $\E[V\mid U]$。

\exercise{AA17}{正态变量平方后的相关性}
$(X,Y)$ 联合标准正态，相关系数为 $\rho\in[-1,1]$。求条件二阶矩 $\E[Y^2\mid X]$，再求两个平方变量的协方差与相关系数。平方后是否还能从相关系数判断原来的相关符号？

\exercise{AA20}{球壳上随机点的距离倒数}
在三维空间中，$X$ 按面积均匀分布于半径 $R>0$ 的球面。固定点 $a$ 到球心的距离为 $r\ge0$，记 $D=|X-a|$。求 $\E(1/D)$；当 $r>0$ 时再求 $\E(1/D^2)$，并判断 $r=R$ 的情形。

\exercise{AA06}{相邻结果首次相同}
独立试验成功概率为 $p\in(0,1)$，记失败概率 $q=1-p$。从零次试验开始，第一次出现连续两次结果相同时停止。求试验总次数的均值，以及因连续两次成功而停止的概率。

\exercise{AA08}{相邻模式出现次数的方差}
独立抛掷正面概率 $p\in(0,1)$ 的硬币 $n\ge2$ 次。统计相邻模式HT的出现次数 $N$，不同起点都计数。求 $\E N$ 与 $\Var N$，并计算 $n=10,p=1/2$ 的结果。

\exercise{AA12}{环游观测的时间相关}
在 $n\ge3$ 个顶点的环上作对称最近邻随机游走，$X_0$ 在全部顶点上均匀分布。令 $Y_t=\cos(2\pi X_t/n)$、$S_T=\sum_{t=0}^{T-1}Y_t$，其中 $T\ge1$。求 $\Cov(Y_0,Y_k)$、$\Var S_T$ 和 $\lim_{T\to\infty}\Var S_T/T$。

\exercise{AA04}{给定尾零数的阶乘反求}
给定非负整数 $K$，寻找所有使 $n!$ 的十进制末尾恰有 $K$ 个零的非负整数 $n$，约定 $0!=1$。设计不用计算阶乘的算法，证明解集若非空必为五个连续整数，并分别求 $K=24$ 与 $K=29$ 的解集。

\exercise{AA05}{二项式系数的末个非零数字}
设计计算 $\binom nk$ 的末个非零十进制数字的算法，其中整数 $0\le k\le n$。不先构造完整的二项式系数，而应分别处理因子2、5及其余单位部分。说明为何直接对模10的分母作除法不合法，并计算 $\binom{20}{7}$。

\exercise{AA11}{随机格点之间的可见性}
从整数网格 $\{0,\ldots,M\}\times\{0,\ldots,N\}$ 中等概率选取两个不同点，其中 $M,N\ge1$。若连接线段内部没有其他整数格点，则称两点互相可见。给出可见概率的有限求和公式，并求 $M=N=2$ 时的值。

\exercise{AA14}{整数三角形的渐近个数}
记 $T(N)$ 为正整数三元组 $x\le y\le z$、$x+y+z=N$ 且 $x+y>z$ 的个数。证明 $T(N)=N^2/48+O(N)$，并列出 $N=12$ 时的全部解。说明直接把有序解个数除以6为何通常不是精确答案。

\exercise{AA15}{接近整数的代数幂}
令 $\alpha=(3+\sqrt5)/2$。构造整数递推，给出对每个 $n\ge1$ 距离 $\alpha^n$ 最近的整数 $L_n$，证明最近整数唯一，并求其误差和奇偶规律。求使误差小于 $10^{-6}$ 的最小 $n$。

\exercise{AA21}{同时约束数位和与数位积}
在十进制正整数中寻找数位积为72且数位和为14的最小数。再把数位和改为18，求新的最小数，并解释为什么“数字1永远可以删除”的说法在这里不成立。

\exercise{AA18}{对数的超越性与比值的无理性}
证明 $\log2$ 是超越数，而 $\log2/\log3$ 是无理数，并说明“两个超越数之比一定超越”为什么错误。这里 $\log$ 表示自然对数。可使用 Hermite--Lindemann 定理：非零代数数 $z$ 的指数 $e^z$ 是超越数。
